\documentclass{article}

% if you need to pass options to natbib, use, e.g.:
%     \PassOptionsToPackage{numbers, compress}{natbib}
% before loading neurips_2026

% The authors should use one of these tracks.
% Before accepting by the NeurIPS conference, select one of the options below.
% 0. "default" for submission
\usepackage{neurips_2026}
% the "default" option is equal to the "main" option, which is used for the Main Track with double-blind reviewing.
% 1. "main" option is used for the Main Track
%  \usepackage[main]{neurips_2026}
% 2. "position" option is used for the Position Paper Track
%  \usepackage[position]{neurips_2026}
% 3. "eandd" option is used for the Evaluations & Datasets Track
 % \usepackage[eandd]{neurips_2026}
 % if you need to opt-in for a single-blind submission in the E&D track:
 %\usepackage[eandd, nonanonymous]{neurips_2026}
% 4. "creativeai" option is used for the Creative AI Track
%  \usepackage[creativeai]{neurips_2026}
% 5. "sglblindworkshop" option is used for the Workshop with single-blind reviewing
 % \usepackage[sglblindworkshop]{neurips_2026}
% 6. "dblblindworkshop" option is used for the Workshop with double-blind reviewing
%  \usepackage[dblblindworkshop]{neurips_2026}

% After being accepted, the authors should add "final" behind the track to compile a camera-ready version.
% 1. Main Track
 % \usepackage[main, final]{neurips_2026}
% 2. Position Paper Track
%  \usepackage[position, final]{neurips_2026}
% 3. Evaluations & Datasets Track
 % \usepackage[eandd, final]{neurips_2026}
% 4. Creative AI Track
%  \usepackage[creativeai, final]{neurips_2026}
% 5. Workshop with single-blind reviewing
%  \usepackage[sglblindworkshop, final]{neurips_2026}
% 6. Workshop with double-blind reviewing
%  \usepackage[dblblindworkshop, final]{neurips_2026}
% Note. For the workshop paper template, both \title{} and \workshoptitle{} are required, with the former indicating the paper title shown in the title and the latter indicating the workshop title displayed in the footnote.
% For workshops (5., 6.), the authors should add the name of the workshop, "\workshoptitle" command is used to set the workshop title.
% \workshoptitle{WORKSHOP TITLE}

% "preprint" option is used for arXiv or other preprint submissions
 % \usepackage[preprint]{neurips_2026}

% to avoid loading the natbib package, add option nonatbib:
%    \usepackage[nonatbib]{neurips_2026}

\usepackage[utf8]{inputenc} % allow utf-8 input
\usepackage[T1]{fontenc}    % use 8-bit T1 fonts
\usepackage{hyperref}       % hyperlinks
\usepackage{url}            % simple URL typesetting
\usepackage{booktabs}       % professional-quality tables
\usepackage{amsfonts}       % blackboard math symbols
\usepackage{nicefrac}       % compact symbols for 1/2, etc.
\usepackage{microtype}      % microtypography
\usepackage{xcolor}         % colors
\usepackage{amsmath}

% Note. For the workshop paper template, both \title{} and \workshoptitle{} are required, with the former indicating the paper title shown in the title and the latter indicating the workshop title displayed in the footnote. 
\title{Beyond Reward Correctness: Distribution-Driven Mechanisms in LLM Reinforcement Learning}


% The \author macro works with any number of authors. There are two commands
% used to separate the names and addresses of multiple authors: \And and \AND.
%
% Using \And between authors leaves it to LaTeX to determine where to break the
% lines. Using \AND forces a line break at that point. So, if LaTeX puts 3 of 4
% authors names on the first line, and the last on the second line, try using
% \AND instead of \And before the third author name.


\author{%
  David S.~Hippocampus\thanks{Use footnote for providing further information
    about author (webpage, alternative address)---\emph{not} for acknowledging
    funding agencies.} \\
  Department of Computer Science\\
  Cranberry-Lemon University\\
  Pittsburgh, PA 15213 \\
  \texttt{hippo@cs.cranberry-lemon.edu} \\
  % examples of more authors
  % \And
  % Coauthor \\
  % Affiliation \\
  % Address \\
  % \texttt{email} \\
  % \AND
  % Coauthor \\
  % Affiliation \\
  % Address \\
  % \texttt{email} \\
  % \And
  % Coauthor \\
  % Affiliation \\
  % Address \\
  % \texttt{email} \\
  % \And
  % Coauthor \\
  % Affiliation \\
  % Address \\
  % \texttt{email} \\
}


\begin{document}


\maketitle


\begin{abstract}
Reinforcement learning (RL) with verifiable rewards (RLVR) has been shown to be an effective post-training approach for Large Language Models (LLMs). Since RL optimizes the model for maximum rewards, it is widely believed that the reward must strongly correlate to task success. However, it does not perfectly explain why the training is robust to noisy self-generated rewards, in self-supervised RL processes. Motivated by recent findings that entropy minimization training includes a distribution-driven process, we hypothesize that such self-supervised RL includes a similar mechanism because it is aiming to minimize the response entropy. We experiment this hypothesis with a design that isolates the success factor with the distribution factor and results show that both of them are effective reward signals, providing evidence to support the existence of the distribution-driven mechanism. Our findings advance current understanding of the LLM RL mechanism, providing explanations to recent findings and informing future research to account for the distribution factor in training.

\end{abstract}

\section{Introduction}

Reinforcement learning (RL) with verifiable rewards (RLVR) has been shown to be an effective post-training approach, especially for recently emerged large language models (LLMs) \cite{guo2025deepseek,zhang2025survey}. Generally, Such approaches reward the model based on success of their action trajectories, and then update model weights to encourage successful behaviors and suppress unsuccessful ones, effectively enhancing the model's problem solving capability \cite{guo2025deepseek,littman2015reinforcement,mnih2015human,silver2016mastering}. In such processes, the reward signal is believed to be the most important factor whose quality directly affects model performance, because it is ultimate goal the model is optimized for \cite{kaelbling1996reinforcement,sutton1998reinforcement}. As such, a great deal of recent work has focused on providing more effective reward signals in RLVR processes \cite{she2025r,wang2024math,yuanfree,guosegment}. 


Recently, it was found that RLVR-style post-training can be done in an self-supervised manner \cite{zuottrl,zhangright,zhangconsistent,agarwalunreasonable}. In particular, the model's own generations were curated as rewards for RL training, such as self-consistency of the answer, e.g, TTRL \cite{zuottrl}, EMPO \cite{zhangright}, internal confidence \cite{agarwalunreasonable}, or self-generated marks \cite{yuan2024self,wu2025meta}. The effectiveness of this training style has been verify by many works, showing comparative performance against training with the ground truth label \cite{zuottrl,zhangright}. It was believed that the key factor behind such processes is the model's self-evaluation \cite{zuottrl,yuan2024self,wu2025meta} and cognitive ability \cite{gandhicognitive}. On the other hand, the quality of such guidance signals are often significantly lower than the ground truth. For example, our experimental results showed that the model can reach a performance comparable to training with the ground truth oracle (often only 1-2\% lower) with only 40\%~70\% reward accuracy. This raises a fundamental question, why such noisy rewards still effectively improve model performance? Why the noise is not harmful?

A possibility of this effect is that there is another mechanism guiding the training in parallel with the rewards. Press et al. revealed that such self-supervised training (i.e., entropy minimization) is not only success-driven, but also a distribution-driven process \cite{press2024entropy}. The model's performance still improves even with 0\% correct labels. A further examination into the training process revealed that it changes the inner representation of the post-training data, mapping the new data to an area close to similar pre-training data. Such vicinity drives the model to act in a way obtained from pre-training, which adapts the model for new tasks.


Given the fact that self-supervised LLM RL training is essentially entropy minimization style as well, we hypothesize that the distribution-driven mechanism also exists and this explains why self-supervised training is robust to noisy rewards. In particular, in Group Relative Policy Optimization (GRPO) style training, the model samples multiple responses for each prompt, and the majority answer is positively rewarded whereas others are negatively rewarded. The diverse responses can be viewed as a probability distribution, and the training aims to enhance the largest group (i.e., class). This is exactly the mechanism proposed by \cite{zuottrl,zhangright}, where Zhang et al. framed it as semantic entropy optimization \cite{zhangright}. 

We experiment this hypothesis with a design that isolates reward from data distribution. Specifically, we pre-sample the responses for the training prompts with the base model to get a landscape of its output distribution, and use it as reference density. Then we train the model with different types of rewards: 1) ground truth (main effects of success and distribution), 2) dynamic self-consistency (main effects of distribution + side effects of rewards), 3) reference density (main effects of distribution + side effects of rewards), 4) correct but sparse (full effects of success), 5) incorrect but dense (full effects of distribution), and 6) random and sparse (no effects of distribution nor success). Results suggest that both the density and reward effects exist, with the effect of success slightly stronger than the effect of distribution. Our work proposes and confirms that distribution-driven processes indeed exist in LLM RL training. Our contributions are three-fold:

1. We propose an hypothesis that LLM RL training is not only success-driven, but also a distribution-driven process.

2. We confirm this hypothesis with experimental results showing that distribution-driven training can achieve comparable model performance as success-driven training.

3. Our findings advance current understanding of the LLM RL mechanism, informing future research to account for the distribution factor in training.


\section{Related Work}

Entropy minimization training (as a type of unsupervised learning) has been found to enhance model performance for a long time \cite{lee2013pseudo,wangtent,niutowards,press2024entropy}. Lee proposed a semi-supervised learning approach where labeled data were trained via standard supervised learning whereas unlabeled data were trained to minimize the model's output entropy \cite{lee2013pseudo}. Analysis showed that entropy minimization training has the clustering effect, reducing within group and distance and increasing between group distance, resulting in clearer separation boundaries. Wang et al. proposed an entropy minimization approach (TENT) for test-time adaptation, where a probability model was trained to minimize the entropy of its own outputs towards the most probable class on new data \cite{wangtent}. Analysis showed that TENT better shapes the model's inner features towards the oracle than other approaches. Niu et al. found that filtering noisy samples and encourage model weights to go to a flat minimum could effectively stabilize entropy minimization test-time adaption \cite{niutowards}. Press et al. digged further into the underlying mechanism of entropy minimization training and discover that such training changes the model's inner representation of new data, mapping them to an area closer to similar training data and thus, similar outputs will be generated \cite{press2024entropy}. More importantly, the model's performance still improves even 100\% labels are incorrect, suggesting that entropy minimization is not totally success driven. To summarize, entropy minimization has been shown to effectively adapt the model for new tasks and it operates mainly through changing the model's inner representation of new data. 

Unsupervised RL for LLM has been explored by a considerable amount of recent studies as RLVR has become one the most attractive areas \cite{zuottrl,zhangright,zhangconsistent,agarwalunreasonable}. Researchers have investigated a variety types of rewards, such as self-consistency \cite{zuottrl,zhangright}, internal confidence \cite{agarwalunreasonable}, and self-generated marks \cite{yuan2024self,wu2025meta}, and cross-model rewards \cite{zhang2025co}. In particular, test-time reinforcement learning (TTRL) positively rewards the majority group (in terms of the final answer) and negatively rewards all other groups, achieving comparable model performance against ground truth rewards. Entropy Minimized Policy Optimization (EMPO) applies a similar rewarding strategy, but weights the rewards relative to group size, effectively enhancing the majority class. Agarwal et al. proposed multiple entropy based objectives for RL training, such as token-level entropy, sequence-level entropy, and token logits \cite{agarwalunreasonable}. Zhang et al. proposed to enhance label quality by referencing to a second view, such as asking superficially different but fundamentally similar questions to the same model, or asking the same question to different model \cite{zhang2025co}. By comparing the answers, it is more likely to obtain the correct label. Put together, prior work has found that enhancing self-consistency can effectively improve model performance. On the other hand, it is widely believed that model performance improves because it is trained towards the true label. Here, we show that training the model towards dense area is another mechanism that can improve model performance. 



\section{Methods}

\subsection{Preliminary: Group Relative Policy Optimization}

Group Relative Policy Optimization (GRPO) \cite{shao2024deepseekmath,guo2025deepseek} is a policy gradient method widely adopted in recent LLM post-training. Given a prompt $q$, GRPO samples a group of $G$ responses $\{o_1, \ldots, o_G\}$ from the current policy $\pi_\theta$ and assigns each response a scalar reward $r_i$ via a reward function. The advantage of each response is computed by normalizing rewards within the group:
\begin{equation}
A_i = \frac{r_i - \text{mean}(r_1, \ldots, r_G)}{\text{std}(r_1, \ldots, r_G)}
\end{equation}
The policy is then updated by maximizing the clipped surrogate objective:
\begin{equation}
\mathcal{J}_{\text{GRPO}} = \mathbb{E}_{q \sim \mathcal{D},\, \{o_i\} \sim \pi_\theta(\cdot|q)} \left[ \frac{1}{G} \sum_{i=1}^{G} \min\left(\rho_i A_i,\; \text{clip}(\rho_i, 1-\epsilon, 1+\epsilon) A_i \right) \right]
\end{equation}
where $\rho_i = \frac{\pi_\theta(o_i|q)}{\pi_{\theta_{\text{old}}}(o_i|q)}$ is the importance sampling ratio between the current and old policy, and $\epsilon$ is the clipping threshold for training stability. Following recent practice \cite{zhangright,zuottrl}, we set the KL penalty coefficient to zero.

In standard RLVR, the reward is computed by comparing the response against a ground-truth answer: $r_i = 1$ if the extracted answer matches the ground truth, and $r_i = 0$ otherwise. This binary reward is the only external signal that enters the optimization --- it determines which responses receive positive advantage and which receive negative advantage.


\subsection{Experimental Logic: Isolating Success from Distribution}

As stated in the introduction, we hypothesize that self-supervised GRPO training contains two separable learning signals: a success factor (reward correctness) and a distribution factor (reward density). To test this, we need experimental conditions that independently vary these two factors while holding the training algorithm constant. The key challenge is that under normal self-supervised training, correctness and density are confounded --- the majority cluster is often correct in well-pretrained models (40\%--70\% accuracy in our setting). To break this confound, we adopt a two-step design:

\paragraph{Step 1: Pre-compute a fixed density reference.} Before training begins, we sample responses ($N=16$) from the frozen base model for every training prompt and cluster them by the final answer, as did by prior studies \cite{zuottrl,zhangright}. The largest cluster's representative answer is stored as the \textit{density target} for that prompt. Because the reference model is frozen, this density landscape is fixed throughout training --- it captures where the model's output distribution is dense at initialization, independent of any learning that occurs during training.

\paragraph{Step 2: Construct reward conditions that factorially vary correctness and density.} Using the pre-computed density targets and ground-truth labels, we partition prompts and assign rewards to create conditions where the two factors are either aligned, isolated, or absent. The critical insight is that for a substantial fraction of prompts (30\%--60\%), the density target disagrees with the ground truth. On these prompts, we can reward one factor while explicitly excluding the other.

\begin{table}[t]
\caption{Evaluation results on different benchmarks. Scores are shown as percentages.}
\label{tab:benchmark_results}
\centering
\small
\setlength{\tabcolsep}{4pt}

\begin{tabular}{lcccccccc}
	oprule
& AMC23 
& AIME24 
& MATH500 
& GSM8K 
& Math Avg.
& \shortstack{MMLU\\Pro}
& \shortstack{GPQA\\Diamond}
& \shortstack{LiveCode\\Bench} \\
\midrule

\multicolumn{9}{l}{\textit{Qwen3-4B-base}} \\
\midrule

base
& 43.4 & 8.3 & 62.6 & 68.3 & 45.6
& 28.6 & 24.4 & 17.4 \\

correct-sparse
& 56.9 & 15.4 & 71.9 & 89.1 & 58.3
& 37.8 & 32.6 & 26.2 \\

ground truth 
& 52.2 & 14.2 & 78.1 & 92.6 & 59.3
& 50.3 & 40.1 & 27.1 \\

incorrect-dense
& 52.5 & 8.8 & 72.0 & 91.0 & 56.1
& 46.5 & 40.8 & 28.2 \\


dense only
& 54.1 & 11.3 & 75.5 & 92.9 & 58.5
& 49.3 & 36.0 & 24.9 \\

self-consistency
& 51.3 & 14.2 & 79.3 & 91.3 & 59.0
& 48.1 & 39.9 & 14.3 \\

\midrule

\multicolumn{9}{l}{\textit{Llama-3.2-3B-Instruct}} \\
\midrule

base
& 17.5 & 6.3 & 41.5 & 71.0 & 34.1
& 25.0 & 20.0 & 5.4 \\

correct-sparse 
& 30.3 & 7.9 & 46.3 & 84.0 & 42.1
& 29.4 & 21.9 & 6.9 \\

ground truth 
& 33.1 & 9.2 & 55.9 & 77.4 & 43.9
& 29.6 & 25.9 & 7.0 \\

incorrect-dense
& 18.4 & 10.8 & 39.3 & 61.9 & 32.6
& 13.6 & 18.9 & 5.8 \\

dense only
& 28.4 & 9.6 & 36.4 & 78.0 & 38.1
& 25.3 & 21.1 & 3.8 \\

self-consistency_1
& 30.0 & 4.2 & 36.9 & 72.1 & 35.8
& 11.5 & 13.0 & 7.5 \\

\bottomrule
\end{tabular}
\end{table}

Papers to be submitted to NeurIPS 2026 must be prepared according to the
instructions presented here. Papers may only be up to {\bf nine} pages long, including figures. \textbf{Papers that exceed the page limit will not be reviewed (or in any other way considered) for presentation at the conference.}
Additional pages \emph{containing acknowledgments, references, checklist, and optional technical appendices} do not count as content pages. 



\section{Submission of papers to NeurIPS 2026}


Please read the instructions below carefully and follow them faithfully.


\subsection{Style}


Papers to be submitted to NeurIPS 2026 must be prepared according to the
instructions presented here. Papers may only be up to {\bf nine} pages long, including figures. \textbf{Papers that exceed the page limit will not be reviewed (or in any other way considered) for presentation at the conference.}
Additional pages \emph{containing acknowledgments, references, checklist, and optional technical appendices} do not count as content pages. 


The margins in 2026 are the same as those in previous years.


Authors are required to use the NeurIPS \LaTeX{} style files obtainable at the NeurIPS website as indicated below. Please make sure you use the current files and not previous versions. Tweaking the style files may be grounds for desk rejection.


\subsection{Retrieval of style files}


The style files for NeurIPS and other conference information are available on the website at
\begin{center}
  \url{https://neurips.cc}.
\end{center}
% The file \verb+neurips_2026.pdf+ contains these instructions and illustrates the various formatting requirements your NeurIPS paper must satisfy.


The only supported style file for NeurIPS 2026 is \verb+neurips_2026.sty+, rewritten for \LaTeXe{}. \textbf{Previous style files for \LaTeX{} 2.09,  Microsoft Word, and RTF are no longer supported.}


The \LaTeX{} style file contains three optional arguments: 
\begin{itemize}
\item \verb+final+, which creates a camera-ready copy, 
\item \verb+preprint+, which creates a preprint for submission to, e.g., arXiv, \item \verb+nonatbib+, which will not load the \verb+natbib+ package for you in case of package clash.
\end{itemize}


\paragraph{Preprint option}
If you wish to post a preprint of your work online, e.g., on arXiv, using the NeurIPS style, please use the \verb+preprint+ option. This will create a nonanonymized version of your work with the text ``Preprint. Work in progress.'' in the footer. This version may be distributed as you see fit, as long as you do not say which conference it was submitted to. Please \textbf{do not} use the \verb+final+ option, which should \textbf{only} be used for papers accepted to NeurIPS.


At submission time, please omit the \verb+final+ and \verb+preprint+ options. This will anonymize your submission and add line numbers to aid review. Please do \emph{not} refer to these line numbers in your paper as they will be removed during generation of camera-ready copies.


The file \verb+neurips_2026.tex+ may be used as a ``shell'' for writing your paper. All you have to do is replace the author, title, abstract, and text of the paper with your own.


The formatting instructions contained in these style files are summarized in Sections \ref{gen_inst}, \ref{headings}, and \ref{others} below.


\section{General formatting instructions}
\label{gen_inst}


The text must be confined within a rectangle 5.5~inches (33~picas) wide and
9~inches (54~picas) long. The left margin is 1.5~inch (9~picas).  Use 10~point
type with a vertical spacing (leading) of 11~points.  Times New Roman is the
preferred typeface throughout, and will be selected for you by default.
Paragraphs are separated by \nicefrac{1}{2}~line space (5.5 points), with no
indentation.


The paper title should be 17~point, initial caps/lower case, bold, centered
between two horizontal rules. The top rule should be 4~points thick and the
bottom rule should be 1~point thick. Allow \nicefrac{1}{4}~inch space above and
below the title to rules. All pages should start at 1~inch (6~picas) from the
top of the page.


For the final version, authors' names are set in boldface, and each name is
centered above the corresponding address. The lead author's name is to be listed
first (left-most), and the co-authors' names (if different address) are set to
follow. If there is only one co-author, list both author and co-author side by
side.


Please pay special attention to the instructions in Section \ref{others}
regarding figures, tables, acknowledgments, and references.

\section{Headings: first level}
\label{headings}


All headings should be lower case (except for first word and proper nouns),
flush left, and bold.


First-level headings should be in 12-point type.


\subsection{Headings: second level}


Second-level headings should be in 10-point type.


\subsubsection{Headings: third level}


Third-level headings should be in 10-point type.


\paragraph{Paragraphs}


There is also a \verb+\paragraph+ command available, which sets the heading in
bold, flush left, and inline with the text, with the heading followed by 1\,em
of space.


\section{Citations, figures, tables, references}
\label{others}


These instructions apply to everyone.


\subsection{Citations within the text}


The \verb+natbib+ package will be loaded for you by default.  Citations may be
author/year or numeric, as long as you maintain internal consistency.  As to the
format of the references themselves, any style is acceptable as long as it is
used consistently.


The documentation for \verb+natbib+ may be found at
\begin{center}
  \url{http://mirrors.ctan.org/macros/latex/contrib/natbib/natnotes.pdf}
\end{center}
Of note is the command \verb+\citet+, which produces citations appropriate for
use in inline text.  For example,
\begin{verbatim}
   \citet{hasselmo} investigated\dots
\end{verbatim}
produces
\begin{quote}
  Hasselmo, et al.\ (1995) investigated\dots
\end{quote}


If you wish to load the \verb+natbib+ package with options, you may add the
following before loading the \verb+neurips_2026+ package:
\begin{verbatim}
   \PassOptionsToPackage{options}{natbib}
\end{verbatim}


If \verb+natbib+ clashes with another package you load, you can add the optional
argument \verb+nonatbib+ when loading the style file:
\begin{verbatim}
   \usepackage[nonatbib]{neurips_2026}
\end{verbatim}


As submission is double blind, refer to your own published work in the third
person. That is, use ``In the previous work of Jones et al.\ [4],'' not ``In our
previous work [4].'' If you cite your other papers that are not widely available
(e.g., a journal paper under review), use anonymous author names in the
citation, e.g., an author of the form ``A.\ Anonymous'' and include a copy of the anonymized paper in the supplementary material.


\subsection{Footnotes}


Footnotes should be used sparingly.  If you do require a footnote, indicate
footnotes with a number\footnote{Sample of the first footnote.} in the
text. Place the footnotes at the bottom of the page on which they appear.
Precede the footnote with a horizontal rule of 2~inches (12~picas).


Note that footnotes are properly typeset \emph{after} punctuation
marks.\footnote{As in this example.}


\subsection{Figures}


\begin{figure}
  \centering
  \fbox{\rule[-.5cm]{0cm}{4cm} \rule[-.5cm]{4cm}{0cm}}
  \caption{Sample figure caption. Explain what the figure shows and add a key take-away message to the caption.}
\end{figure}


All artwork must be neat, clean, and legible. Lines should be dark enough for
 reproduction purposes. The figure number and caption always appear after the
figure. Place one line space before the figure caption and one line space after
the figure. The figure caption should be lower case (except for the first word and proper nouns); figures are numbered consecutively.


You may use color figures.  However, it is best for the figure captions and the
paper body to be legible if the paper is printed in either black/white or in
color.


\subsection{Tables}


All tables must be centered, neat, clean, and legible.  The table number and
title always appear before the table.  See Table~\ref{sample-table}.


Place one line space before the table title, one line space after the
table title, and one line space after the table. The table title must
be lower case (except for the first word and proper nouns); tables are
numbered consecutively.


Note that publication-quality tables \emph{do not contain vertical rules}. We
strongly suggest the use of the \verb+booktabs+ package, which allows for
typesetting high-quality, professional tables:
\begin{center}
  \url{https://www.ctan.org/pkg/booktabs}
\end{center}
This package was used to typeset Table~\ref{sample-table}.


\begin{table}
  \caption{Sample table caption. Explain what the table shows and add a key take-away message to the caption.}
  \label{sample-table}
  \centering
  \begin{tabular}{lll}
    \toprule
    \multicolumn{2}{c}{Part}                   \\
    \cmidrule(r){1-2}
    Name     & Description     & Size ($\mu$m) \\
    \midrule
    Dendrite & Input terminal  & $\approx$100     \\
    Axon     & Output terminal & $\approx$10      \\
    Soma     & Cell body       & up to $10^6$  \\
    \bottomrule
  \end{tabular}
\end{table}

\subsection{Math}
Note that display math in bare TeX commands will not create correct line numbers for submission. Please use LaTeX (or AMSTeX) commands for unnumbered display math. (You really shouldn't be using \$\$ anyway; see \url{https://tex.stackexchange.com/questions/503/why-is-preferable-to} and \url{https://tex.stackexchange.com/questions/40492/what-are-the-differences-between-align-equation-and-displaymath} for more information.)

\subsection{Final instructions}

Do not change any aspects of the formatting parameters in the style files.  In
particular, do not modify the width or length of the rectangle the text should
fit into, and do not change font sizes. Please note that pages should be
numbered.


\section{Preparing PDF files}


Please prepare submission files with paper size ``US Letter,'' and not, for
example, ``A4.''


Fonts were the main cause of problems in the past years. Your PDF file must only
contain Type 1 or Embedded TrueType fonts. Here are a few instructions to
achieve this.


\begin{itemize}


\item You should directly generate PDF files using \verb+pdflatex+.


\item You can check which fonts a PDF files uses.  In Acrobat Reader, select the
  menu Files$>$Document Properties$>$Fonts and select Show All Fonts. You can
  also use the program \verb+pdffonts+ which comes with \verb+xpdf+ and is
  available out-of-the-box on most Linux machines.


\item \verb+xfig+ ``patterned'' shapes are implemented with bitmap fonts.  Use
  "solid" shapes instead.


\item The \verb+\bbold+ package almost always uses bitmap fonts.  You should use
  the equivalent AMS Fonts:
\begin{verbatim}
   \usepackage{amsfonts}
\end{verbatim}
followed by, e.g., \verb+\mathbb{R}+, \verb+\mathbb{N}+, or \verb+\mathbb{C}+
for $\mathbb{R}$, $\mathbb{N}$ or $\mathbb{C}$.  You can also use the following
workaround for reals, natural and complex:
\begin{verbatim}
   \newcommand{\RR}{I\!\!R} %real numbers
   \newcommand{\Nat}{I\!\!N} %natural numbers
   \newcommand{\CC}{I\!\!\!\!C} %complex numbers
\end{verbatim}
Note that \verb+amsfonts+ is automatically loaded by the \verb+amssymb+ package.


\end{itemize}


If your file contains type 3 fonts or non embedded TrueType fonts, we will ask
you to fix it.


\subsection{Margins in \LaTeX{}}


Most of the margin problems come from figures positioned by hand using
\verb+\special+ or other commands. We suggest using the command
\verb+\includegraphics+ from the \verb+graphicx+ package. Always specify the
figure width as a multiple of the line width as in the example below:
\begin{verbatim}
   \usepackage[pdftex]{graphicx} ...
   \includegraphics[width=0.8\linewidth]{myfile.pdf}
\end{verbatim}
See Section 4.4 in the graphics bundle documentation
(\url{http://mirrors.ctan.org/macros/latex/required/graphics/grfguide.pdf})


A number of width problems arise when \LaTeX{} cannot properly hyphenate a
line. Please give LaTeX hyphenation hints using the \verb+\-+ command when
necessary.

\begin{ack}
Use unnumbered first level headings for the acknowledgments. All acknowledgments
go at the end of the paper before the list of references. Moreover, you are required to declare
funding (financial activities supporting the submitted work) and competing interests (related financial activities outside the submitted work).
More information about this disclosure can be found at: \url{https://neurips.cc/Conferences/2026/PaperInformation/FundingDisclosure}.


Do {\bf not} include this section in the anonymized submission, only in the final paper. You can use the \texttt{ack} environment provided in the style file to automatically hide this section in the anonymized submission.
\end{ack}


\bibliographystyle{plainnat}
\bibliography{references}

% \section*{References}


% References follow the acknowledgments in the camera-ready paper. Use unnumbered first-level heading for
% the references. Any choice of citation style is acceptable as long as you are
% consistent. It is permissible to reduce the font size to \verb+small+ (9 point)
% when listing the references.
% Note that the Reference section does not count towards the page limit.
% \medskip


% {
% \small


% [1] Alexander, J.A.\ \& Mozer, M.C.\ (1995) Template-based algorithms for
% connectionist rule extraction. In G.\ Tesauro, D.S.\ Touretzky and T.K.\ Leen
% (eds.), {\it Advances in Neural Information Processing Systems 7},
% pp.\ 609--616. Cambridge, MA: MIT Press.


% [2] Bower, J.M.\ \& Beeman, D.\ (1995) {\it The Book of GENESIS: Exploring
%   Realistic Neural Models with the GEneral NEural SImulation System.}  New York:
% TELOS/Springer--Verlag.


% [3] Hasselmo, M.E., Schnell, E.\ \& Barkai, E.\ (1995) Dynamics of learning and
% recall at excitatory recurrent synapses and cholinergic modulation in rat
% hippocampal region CA3. {\it Journal of Neuroscience} {\bf 15}(7):5249-5262.
% }


%%%%%%%%%%%%%%%%%%%%%%%%%%%%%%%%%%%%%%%%%%%%%%%%%%%%%%%%%%%%

\appendix

\section{Technical appendices and supplementary material}
Technical appendices with additional results, figures, graphs, and proofs may be submitted with the paper submission before the full submission deadline (see above). You can upload a ZIP file for videos or code, but do not upload a separate PDF file for the appendix. There is no page limit for the technical appendices. 

Note: Think of the appendix as ``optional reading'' for reviewers. The paper must be able to stand alone without the appendix; for example, adding critical experiments that support the main claims to an appendix is inappropriate. 

%%%%%%%%%%%%%%%%%%%%%%%%%%%%%%%%%%%%%%%%%%%%%%%%%%%%%%%%%%%%

\newpage
\input{checklist.tex}


\end{document}